\documentclass[11pt,a4paper]{article}
% Preâmbulo comum. Não duplicar pacotes nos módulos.
\usepackage[utf8]{inputenc}
\usepackage[T1]{fontenc}
\usepackage[brazil]{babel}
\usepackage{amsmath,amssymb,amsthm,mathtools}
\usepackage{geometry}
\geometry{a4paper,margin=2.3cm}
\usepackage{microtype}
\usepackage{enumitem}
\usepackage{graphicx}
\usepackage{xcolor}
\usepackage{booktabs}
\usepackage{hyperref}
\hypersetup{colorlinks=true,linkcolor=blue!40!black,urlcolor=blue!50!black}

\setcounter{tocdepth}{2}
\numberwithin{equation}{section}

\theoremstyle{plain}
\newtheorem{teorema}{Teorema}[section]
\newtheorem{proposicao}[teorema]{Proposição}
\newtheorem{lema}[teorema]{Lema}
\newtheorem{corolario}[teorema]{Corolário}

\theoremstyle{definition}
\newtheorem{definicao}[teorema]{Definição}
\newtheorem{exemplo}[teorema]{Exemplo}
\newtheorem{exercicio}{Exercício}[section]

\theoremstyle{remark}
\newtheorem{observacao}[teorema]{Observação}
\newtheorem{contraexemplo}[teorema]{Contraexemplo}
\newtheorem{intuicao}[teorema]{Intuição}

\newcommand{\R}{\mathbb{R}}
\newcommand{\N}{\mathbb{N}}
\newcommand{\Z}{\mathbb{Z}}
\newcommand{\Q}{\mathbb{Q}}
\DeclareRobustCommand{\eps}{\varepsilon}
\DeclareRobustCommand{\vet}[1]{\vec{#1}}
\DeclareRobustCommand{\dif}{\mathrm{d}}
\DeclareRobustCommand{\abs}[1]{\left\lvert #1 \right\rvert}
\DeclareRobustCommand{\norm}[1]{\left\lVert #1 \right\rVert}
\DeclareMathOperator{\grad}{grad}
\DeclareMathOperator{\Div}{div}
\DeclareMathOperator{\rot}{rot}

\newcommand{\base}[1]{\par\medskip\noindent\textbf{Base para relembrar.} #1\par\medskip}
\newcommand{\cuidado}[1]{\par\medskip\noindent\textbf{Quando a hipótese falha.} #1\par\medskip}

\newcommand{\figura}[2]{%
  \begin{figure}[htbp]
  \centering
  \IfFileExists{#1}{%
    \includegraphics[width=0.84\textwidth]{#1}%
  }{%
    \fbox{\parbox{0.8\textwidth}{\centering Figura ainda não gerada:\\ \texttt{#1}}}%
  }
  \caption{#2}
  \end{figure}%
}

\newenvironment{esqueleto}{%
  \par\medskip\noindent\textbf{Esqueleto da demonstração.}\par\nopagebreak
  \begin{enumerate}[leftmargin=1.4em]
}{\end{enumerate}\medskip}

\title{Cálculo III: funções de várias variáveis}
\author{}
\date{}
\begin{document}
\maketitle
\tableofcontents

\section{Funções, domínio, imagem, gráfico e níveis}

\base{Em uma variável, o domínio de $f$ é o maior subconjunto de $\R$ em que a fórmula faz sentido, e a imagem é o conjunto dos valores efetivamente atingidos. A mesma distinção vale em $\R^n$. A norma euclidiana de $\vet{v}=(v_1,\ldots,v_n)$ é $\norm{\vet{v}}=\sqrt{v_1^2+\cdots+v_n^2}$.}

Uma função escalar de várias variáveis é uma regra $f\colon D\subset\R^n\to\R$. Um campo vetorial é uma regra $\vet{F}\colon D\subset\R^n\to\R^m$, com $m\ge 2$. Escrever $\vet{F}=(F_1,\ldots,F_m)$ é o mesmo que dar $m$ funções escalares, as componentes, todas com o mesmo domínio.

O \emph{domínio} é o conjunto dos pontos em que a expressão está definida nos reais: não se extrai raiz quadrada de número negativo, não se divide por zero, e o logaritmo só recebe argumento positivo. A \emph{imagem} é $\{f(\vet{x}):\vet{x}\in D\}$. Domínio grande não decide a imagem: a função pode deixar de fora muitos reais.

\begin{exemplo}
Seja $f(x,y)=\sqrt{4-x^2-y^2}$. A raiz está definida exatamente quando $4-x^2-y^2\ge 0$, isto é, no disco fechado $x^2+y^2\le 4$. O radicando percorre $[0,4]$: vale $4$ na origem e $0$ na circunferência $x^2+y^2=4$. Logo a imagem de $f$ é $[0,2]$.

O campo $\vet{F}(x,y)=(-y,x)$ tem domínio $\R^2$. A imagem também é $\R^2$: dado $(a,b)$, o ponto $(x,y)=(b,-a)$ satisfaz $\vet{F}(b,-a)=(a,b)$.
\end{exemplo}

O \emph{gráfico} de $f\colon D\subset\R^2\to\R$ é o conjunto
\[
\{(x,y,z)\in\R^3:(x,y)\in D,\ z=f(x,y)\}.
\]
É uma superfície em $\R^3$, no sentido geométrico do conjunto de pontos. Não é o mesmo objeto que uma curva de nível.

\begin{definicao}
Seja $f\colon D\subset\R^2\to\R$ e $c\in\R$. A \emph{curva de nível} $c$ é o conjunto $\{(x,y)\in D:f(x,y)=c\}$. Se $f\colon D\subset\R^3\to\R$, o conjunto $f(x,y,z)=c$ é uma \emph{superfície de nível}.
\end{definicao}

\begin{exemplo}
Para $f(x,y)=x^2+y^2$, a curva de nível $c$ é vazia se $c<0$, é só a origem se $c=0$, e é a circunferência de raio $\sqrt{c}$ se $c>0$.

Para $g(x,y)=x^2-y^2$, a curva de nível $c>0$ é a hipérbole que abre ao longo do eixo $x$; se $c<0$, a hipérbole abre ao longo do eixo $y$; se $c=0$, são as retas $y=\pm x$.
\end{exemplo}

\figura{../figuras/curvas_nivel.png}{Curvas de nível de $x^2+y^2$ e de $x^2-y^2$.}

\begin{intuicao}
O gráfico guarda o valor de $f$ na altura $z$. A curva de nível esquece a altura e marca, no domínio, onde esse valor é constante. Um mapa topográfico é um desenho de curvas de nível da altitude. O gráfico seria a montanha inteira.
\end{intuicao}

\section{Aberto, fechado, limitado, compacto}

\begin{definicao}
A \emph{bola aberta} de centro $\vet{a}\in\R^n$ e raio $r>0$ é
\[
B_r(\vet{a})=\{\vet{x}\in\R^n:\norm{\vet{x}-\vet{a}}<r\}.
\]
A bola fechada troca $<$ por $\le$.
\end{definicao}

\begin{definicao}
Seja $A\subset\R^n$ e $\vet{p}\in\R^n$.
\begin{enumerate}[leftmargin=1.4em]
\item $\vet{p}$ é \emph{ponto interior} de $A$ quando existe $r>0$ com $B_r(\vet{p})\subset A$. O \emph{interior} de $A$ é o conjunto desses pontos.
\item $\vet{p}$ é \emph{ponto de fronteira} de $A$ quando toda bola $B_r(\vet{p})$ contém pelo menos um ponto de $A$ e pelo menos um ponto de $\R^n\setminus A$. A \emph{fronteira} é o conjunto desses pontos.
\item $A$ é \emph{aberto} quando todo ponto de $A$ é interior.
\item $A$ é \emph{limitado} quando existe $M>0$ tal que $\norm{\vet{x}}\le M$ para todo $\vet{x}\in A$.
\item $A$ é \emph{compacto} quando é fechado e limitado. O significado de fechado está no teorema seguinte.
\end{enumerate}
\end{definicao}

Todo ponto de $\R^n$ cai em exatamente uma de três classes em relação a $A$: é interior a $A$, é interior ao complementar, ou é ponto de fronteira. De fato, ou existe uma bola inteira dentro de $A$, ou existe uma bola inteira dentro do complementar, ou toda bola encontra os dois conjuntos.

\begin{teorema}
Para $A\subset\R^n$, as afirmações seguintes são equivalentes.
\begin{enumerate}[leftmargin=1.4em]
\item O complementar $\R^n\setminus A$ é aberto.
\item $A$ contém todos os seus pontos de fronteira.
\end{enumerate}
Um conjunto que satisfaz qualquer uma das duas chama-se \emph{fechado}.
\end{teorema}

\begin{proof}
Suponha o complementar aberto e seja $\vet{p}$ um ponto de fronteira. Se $\vet{p}$ estivesse fora de $A$, seria um ponto do complementar, e haveria uma bola em torno de $\vet{p}$ inteiramente fora de $A$. Essa bola não encontra $A$, e $\vet{p}$ não seria de fronteira. Logo $\vet{p}\in A$.

Reciprocamente, suponha que $A$ contém a fronteira. Tome $\vet{p}\notin A$. Esse ponto não é interior a $A$ e não é de fronteira. Pela tricotomia, é interior ao complementar: alguma bola em torno de $\vet{p}$ fica fora de $A$. O complementar é aberto.
\end{proof}

A bola aberta é aberta: se $\norm{\vet{x}-\vet{a}}<r$ e $\rho=r-\norm{\vet{x}-\vet{a}}>0$, a bola $B_\rho(\vet{x})$ fica dentro de $B_r(\vet{a})$, porque a desigualdade triangular dá, para $\norm{\vet{y}-\vet{x}}<\rho$,
\[
\norm{\vet{y}-\vet{a}}\le\norm{\vet{y}-\vet{x}}+\norm{\vet{x}-\vet{a}}<\rho+\norm{\vet{x}-\vet{a}}=r.
\]
A bola fechada é fechada e limitada, logo compacta. Uma reta em $\R^2$ é fechada e não é limitada.

\begin{teorema}[Heine--Borel]
Seja $K\subset\R^n$. As afirmações seguintes são equivalentes.
\begin{enumerate}[leftmargin=1.4em]
\item $K$ é fechado e limitado.
\item Toda cobertura de $K$ por conjuntos abertos de $\R^n$ admite uma subcobertura finita: se $K\subset\bigcup_{i\in I} U_i$ com cada $U_i$ aberto, então existem $i_1,\ldots,i_N$ tais que $K\subset U_{i_1}\cup\cdots\cup U_{i_N}$.
\end{enumerate}
\end{teorema}

Não provamos essa equivalência. Daqui em diante, compacto quer dizer fechado e limitado, que é a forma operacional do teorema em $\R^n$. A segunda formulação é o que se usa quando se fala de cobertura; ela não será necessária nas contas deste texto. O ganho concreto do teorema, para nós, é o teorema de Weierstrass na nota de otimização: função contínua em compacto atinge máximo e mínimo.

\section{Limite}

\begin{definicao}
Seja $f$ definida numa bola $B_r(a,b)$, exceto possivelmente no próprio ponto $(a,b)$. Escrevemos
\[
\lim_{(x,y)\to(a,b)} f(x,y)=L
\]
quando, para todo $\eps>0$, existe $\delta>0$ tal que
\[
0<\norm{(x,y)-(a,b)}<\delta\quad\Rightarrow\quad\abs{f(x,y)-L}<\eps.
\]
A norma é a euclidiana. O valor $f(a,b)$, mesmo que esteja definido, não entra na definição: a desigualdade pede distância estritamente positiva até $(a,b)$.
\end{definicao}

\begin{intuicao}
O ponto $(x,y)$ pode chegar a $(a,b)$ por qualquer caminho dentro do disco. O número $L$ só é o limite se $f$ se aproximar de $L$ em todos esses modos de chegada ao mesmo tempo. Por isso uma conta ao longo de uma reta só é um teste, não uma demonstração.
\end{intuicao}

\subsection{Normas equivalentes e a norma do máximo}

Em $\R^n$, duas normas $\norm{\cdot}_{\alpha}$ e $\norm{\cdot}_{\beta}$ são \emph{equivalentes} quando existem constantes $c,C>0$ tais que
\[
c\,\norm{\vet{v}}_{\alpha}\le\norm{\vet{v}}_{\beta}\le C\,\norm{\vet{v}}_{\alpha}
\]
para todo $\vet{v}$. É um teorema de álgebra linear que todas as normas em $\R^n$ são equivalentes. Não provamos o caso geral. A consequência que nos importa é imediata a partir das constantes: as bolas de uma norma contêm bolas da outra, com o raio ajustado por $c$ e $C$. Portanto a existência do limite, e o valor do limite, não dependem de qual norma equivalente se usa na definição. Em particular, tanto faz usar a norma euclidiana ou a norma do máximo
\[
\norm{(x,y)}_{\infty}=\max\{\abs{x},\abs{y}\}.
\]
O que provamos agora é a comparação explícita no caso $n=2$, que já basta para trocar uma pela outra nas contas de $\eps$ e $\delta$.

\begin{teorema}
Para todo $\vet{v}=(x,y)\in\R^2$,
\[
\norm{\vet{v}}_{\infty}\le\norm{\vet{v}}_{2}\le\sqrt{2}\,\norm{\vet{v}}_{\infty}.
\]
\end{teorema}

\begin{proof}
Seja $M=\norm{\vet{v}}_{\infty}=\max\{\abs{x},\abs{y}\}$. Então $x^2\le M^2$ e $y^2\le M^2$, logo
\[
\norm{\vet{v}}_{2}^{2}=x^2+y^2\le 2M^2,
\]
e, como os dois lados são não negativos, $\norm{\vet{v}}_{2}\le\sqrt{2}\,M$.

Por outro lado, $M^2$ é $x^2$ ou $y^2$, conforme o maior módulo. Nos dois casos $M^2\le x^2+y^2$, logo $M\le\norm{\vet{v}}_{2}$.
\end{proof}

Com essa conta, os limites coincidem. Suponha $\lim f=L$ no sentido euclidiano, e seja $\eps>0$. Existe $\delta>0$ tal que $0<\norm{\vet{p}-\vet{p}_0}_2<\delta$ implica $\abs{f-L}<\eps$. Se $0<\norm{\vet{p}-\vet{p}_0}_{\infty}<\delta/\sqrt{2}$, a desigualdade da direita dá
\[
\norm{\vet{p}-\vet{p}_0}_{2}\le\sqrt{2}\,\norm{\vet{p}-\vet{p}_0}_{\infty}<\delta,
\]
e portanto $\abs{f-L}<\eps$. O limite na norma do máximo também é $L$.

No outro sentido, $\norm{\cdot}_{\infty}\le\norm{\cdot}_{2}$, então uma bola euclidiana de raio $\delta$ já está dentro da bola do máximo de raio $\delta$. O mesmo $\delta$ serve, sem encolhê-lo.

\subsection{Unicidade}

\begin{teorema}[Unicidade do limite]
Se $\lim_{(x,y)\to(a,b)} f(x,y)=L$ e $\lim_{(x,y)\to(a,b)} f(x,y)=M$, então $L=M$.
\end{teorema}

\begin{proof}
Suponha $L\neq M$ e ponha $\eps=\abs{L-M}/2>0$. Existem $\delta_1,\delta_2>0$ tais que
\[
0<\norm{(x,y)-(a,b)}<\delta_1\Rightarrow\abs{f(x,y)-L}<\eps,
\]
\[
0<\norm{(x,y)-(a,b)}<\delta_2\Rightarrow\abs{f(x,y)-M}<\eps.
\]
Seja $\delta=\min\{\delta_1,\delta_2,r\}$, onde $r>0$ é um raio em que $f$ está definida fora de $(a,b)$. A bola furada de raio $\delta$ não é vazia: o ponto $(a+\delta/2,\,b)$ está nela e no domínio. Para esse ponto,
\[
\abs{L-M}\le\abs{L-f(x,y)}+\abs{f(x,y)-M}<\eps+\eps=\abs{L-M},
\]
o que é impossível. Logo $L=M$.
\end{proof}

\subsection{Limite ao longo de curvas}

\begin{teorema}
Suponha $\lim_{\vet{p}\to\vet{p}_0} f(\vet{p})=L$, com $\vet{p}_0=(a,b)$. Seja $\gamma$ definida numa vizinhança de $t_0\in\R$, com valores em $\R^2$, tal que $\gamma(t)\to\vet{p}_0$ quando $t\to t_0$ e $\gamma(t)\neq\vet{p}_0$ para todo $t\neq t_0$ suficientemente próximo de $t_0$. Então
\[
\lim_{t\to t_0} f(\gamma(t))=L.
\]
Em particular, isso vale para toda curva contínua com $\gamma(t_0)=\vet{p}_0$ que não retorna a $\vet{p}_0$ noutros instantes próximos de $t_0$.
\end{teorema}

\begin{proof}
Seja $\eps>0$. Pela definição do limite em $\R^2$, existe $\delta>0$ tal que $0<\norm{\vet{p}-\vet{p}_0}<\delta$ implica $\abs{f(\vet{p})-L}<\eps$. Como $\gamma(t)\to\vet{p}_0$, existe $\eta>0$ tal que $0<\abs{t-t_0}<\eta$ implica $\norm{\gamma(t)-\vet{p}_0}<\delta$. A hipótese sobre $\gamma$ permite escolher $\eta$ de modo que, além disso, $\gamma(t)\neq\vet{p}_0$. Para esses $t$,
\[
0<\norm{\gamma(t)-\vet{p}_0}<\delta,
\]
logo $\abs{f(\gamma(t))-L}<\eps$.
\end{proof}

A contrapositiva é a ferramenta do dia a dia. Se existirem dois caminhos, ambos chegando a $(a,b)$ sem permanecer em $(a,b)$, ao longo dos quais os limites de $f$ existem e são diferentes, então $\lim_{(x,y)\to(a,b)} f(x,y)$ não existe. Se ao longo de um único caminho o limite não existe, o limite em $\R^2$ também não existe.

\cuidado{Limites iguais ao longo de todas as retas que passam por $(a,b)$ não provam que o limite existe. A reta é uma família muito estreita de caminhos. O teorema acima exige o controle de toda curva que chega ao ponto, e a definição exige o disco inteiro. As duas contas seguintes separam as duas falhas: na primeira, as próprias retas já discordam; na segunda, todas as retas concordam e mesmo assim o limite não existe.}

\begin{contraexemplo}
Seja
\[
f(x,y)=\frac{xy}{x^2+y^2},\qquad (x,y)\neq(0,0).
\]
O valor $f(0,0)$ é irrelevante para o limite. Ao longo da reta $y=mx$, com $m$ fixo e $x\neq 0$,
\begin{align*}
f(x,mx)
&=\frac{x\cdot(mx)}{x^2+(mx)^2}
=\frac{m x^2}{x^2(1+m^2)}
=\frac{m}{1+m^2}.
\end{align*}
O quociente não depende de $x$. Portanto
\[
\lim_{x\to 0} f(x,mx)=\frac{m}{1+m^2}.
\]
Esse número depende de $m$. Para $m=0$, a reta é o eixo $x$ e o limite é $0$. Para $m=1$, a reta é $y=x$ e o limite é $1/2$. Como dois caminhos dão limites diferentes, $\lim_{(x,y)\to(0,0)} f(x,y)$ não existe.

A função $m\mapsto m/(1+m^2)$ percorre $[-1/2,1/2]$: a derivada $(1-m^2)/(1+m^2)^2$ se anula em $m=\pm 1$, e os valores são $\pm 1/2$. Toda reta produz um limite nesse intervalo, e retas diferentes podem produzir limites diferentes.
\end{contraexemplo}

\begin{contraexemplo}
Seja
\[
f(x,y)=\frac{x y^2}{x^2+y^4},\qquad (x,y)\neq(0,0).
\]
De novo, $f(0,0)$ não entra no limite. Ao longo de $y=mx$,
\begin{align*}
f(x,mx)
&=\frac{x\cdot m^2 x^2}{x^2+m^4 x^4}
=\frac{m^2 x^3}{x^2(1+m^4 x^2)}
=\frac{m^2 x}{1+m^4 x^2}.
\end{align*}
Quando $x\to 0$, o numerador tende a $0$ e o denominador tende a $1$, logo o limite ao longo da reta é $0$. Ao longo do eixo $y$, que é a reta vertical $x=0$, o numerador é identicamente zero, então $f(0,y)=0$ e o limite também é $0$. Toda reta que passa pela origem dá limite $0$.

Ao longo da parábola $x=y^2$, com $y\neq 0$,
\[
f(y^2,y)=\frac{y^2\cdot y^2}{(y^2)^2+y^4}=\frac{y^4}{y^4+y^4}=\frac{1}{2}.
\]
O limite ao longo desse caminho é $1/2$, não é $0$. Portanto o limite em $\R^2$ não existe.

A mesma conta com a família $x=c y^2$, $c$ fixo e $y\neq 0$, dá
\[
f(c y^2,y)=\frac{c y^4}{c^2 y^4+y^4}=\frac{c}{c^2+1},
\]
que depende de $c$. A parábola $x=y^2$ é o caso $c=1$. Uma família de parábolas já produz limites diferentes; uma só, diferente do valor das retas, basta para concluir que o limite não existe.
\end{contraexemplo}

\section{Como provar que o limite existe}

Para provar que o limite é $L$, é preciso uma estimativa que valha no disco inteiro, não só numa curva. O caminho padrão é achar uma função $\varphi(r)$, com $\varphi(r)\to 0$ quando $r\to 0$, tal que
\[
\abs{f(x,y)-L}\le \varphi\bigl(\norm{(x,y)-(a,b)}\bigr)
\]
perto do ponto, fora dele. Dado $\eps>0$, escolhe-se $\delta$ a partir de $\varphi$ apenas.

\begin{exemplo}
\[
\lim_{(x,y)\to(0,0)}\frac{x^2 y}{x^2+y^2}=0.
\]
Para $(x,y)\neq(0,0)$,
\[
\abs{\frac{x^2 y}{x^2+y^2}}
=\abs{y}\cdot\frac{x^2}{x^2+y^2}.
\]
Como $x^2\le x^2+y^2$, a fração é no máximo $1$, e
\[
\abs{\frac{x^2 y}{x^2+y^2}}\le\abs{y}.
\]
Além disso $\abs{y}\le\sqrt{x^2+y^2}$. Dado $\eps>0$, tome $\delta=\eps$. Se $0<\sqrt{x^2+y^2}<\delta$, então
\[
\abs{\frac{x^2 y}{x^2+y^2}}\le\abs{y}\le\sqrt{x^2+y^2}<\eps.
\]
A cota $\abs{y}$ já é uma função da distância, porque $\abs{y}$ é ainda menor que a distância. Não há folga escondida: quando $y=0$ e $x\neq 0$ a fração $x^2/(x^2+y^2)$ vale $1$, então $\le\abs{y}$ é afiada ao longo do eixo $x$, onde os dois lados são zero.
\end{exemplo}

\subsection{Polares}

A mudança
\[
x=r\cos\theta,\qquad y=r\sin\theta,\qquad r=\sqrt{x^2+y^2}\ge 0,
\]
transforma a distância à origem na variável $r$. O limite $(x,y)\to(0,0)$ vira $r\to 0$. Ele existe e vale $L$ quando a expressão em $r$ e $\theta$ tende a $L$ com $\delta$ independente de $\theta$: para todo $\eps>0$ existe $\delta>0$ tal que $0<r<\delta$ implica $\abs{f(r\cos\theta,r\sin\theta)-L}<\eps$, qualquer que seja $\theta$.

Se, depois da substituição, aparece um fator $r$ multiplicando uma função limitada em $\theta$, o limite é zero. Se o resultado ainda depende de $\theta$ no limite $r\to 0$, o limite não existe. Polares não são uma definição nova; são uma conta dentro da definição euclidiana.

\begin{exemplo}
No exemplo anterior, com $r>0$,
\[
\frac{x^2 y}{x^2+y^2}
=\frac{r^2\cos^2\theta\cdot r\sin\theta}{r^2}
=r\cos^2\theta\sin\theta.
\]
Logo o módulo é no máximo $r$, porque $\abs{\cos^2\theta\sin\theta}\le 1$. O majorante $r$ não depende de $\theta$ e tende a $0$. O mesmo $\delta=\eps$ funciona para todo ângulo.
\end{exemplo}

\begin{exemplo}
Para $f(x,y)=xy/(x^2+y^2)$, a mesma substituição dá
\[
f(r\cos\theta,r\sin\theta)=\cos\theta\sin\theta=\frac{1}{2}\sin(2\theta),
\]
que não depende de $r$ e depende de $\theta$. Quando $r\to 0$ a expressão fica parada em valores diferentes para ângulos diferentes. Isso recupera, de uma vez, a família de retas: a reta de inclinação $m$ corresponde a um ângulo com $\tan\theta=m$, e $\frac{1}{2}\sin(2\theta)$ reproduz $m/(1+m^2)$.
\end{exemplo}

\section{Continuidade}

\begin{definicao}
Seja $f$ definida num conjunto $D\subset\R^2$ e $(a,b)\in D$. Dizemos que $f$ é \emph{contínua} em $(a,b)$ quando
\[
\lim_{(x,y)\to(a,b)} f(x,y)=f(a,b).
\]
A função é contínua em $D$ quando é contínua em cada ponto de $D$. Nos pontos interiores a definição é a de cima. Num ponto de fronteira de $D$, o limite é entendido com $(x,y)\in D$: a desigualdade $\abs{f-L}<\eps$ só é exigida para pontos do domínio.
\end{definicao}

O limite, quando existe, ignora o valor no ponto. A continuidade pede que esse valor coincida com o limite. No primeiro contraexemplo, qualquer escolha de $f(0,0)$ deixa a função descontínua na origem, porque o limite não existe.

\begin{teorema}[Álgebra de limites]
Se $\lim f=L$ e $\lim g=M$ no mesmo ponto $(a,b)$, então
\[
\lim(f+g)=L+M,\qquad \lim(fg)=LM.
\]
Se $M\neq 0$, então $\lim(f/g)=L/M$.
\end{teorema}

\begin{proof}
Provamos a soma; o produto fica como exercício, e o quociente sai do produto com o inverso. Seja $\eps>0$. Existem $\delta_1$ e $\delta_2$ tais que $0<\norm{(x,y)-(a,b)}<\delta_1$ implica $\abs{f-L}<\eps/2$, e o análogo para $g$ e $M$ com $\delta_2$. Para $\delta=\min\{\delta_1,\delta_2\}$,
\[
\abs{(f+g)-(L+M)}\le\abs{f-L}+\abs{g-M}<\eps.
\]
\end{proof}

\begin{corolario}
Toda função polinomial $p\colon\R^n\to\R$ é contínua em $\R^n$. Toda função racional, quociente de polinômios, é contínua em cada ponto em que o denominador não se anula.
\end{corolario}

\begin{proof}
As projeções $(x,y)\mapsto x$ e $(x,y)\mapsto y$ são contínuas: $\abs{x-a}\le\norm{(x,y)-(a,b)}$, então $\delta=\eps$ serve para a primeira, e o mesmo para a segunda. Constantes são contínuas. Soma e produto de funções contínuas são contínuas, pela álgebra de limites aplicada a $f-L$ no ponto. Um polinômio é obtido de projeções e constantes por um número finito de somas e produtos, logo é contínuo.

Se $p$ e $q$ são polinômios e $q(a,b)\neq 0$, então $p/q$ é quociente de contínuas com denominador não nulo no ponto, e a álgebra de limites dá a continuidade em $(a,b)$. O mesmo argumento, coordenada a coordenada, vale em $\R^n$.
\end{proof}

Fora dos zeros do denominador, não há o que discutir: a expressão racional é composição algébrica de funções contínuas. O trabalho de verdade começa onde o denominador se anula e alguém redefine a função naquele ponto, ou onde aparece uma raiz, um logaritmo, ou uma fórmula por casos.

\section{Exercícios}

\begin{exercicio}[{[Básico --- domínio e imagem]}]
\begin{enumerate}[label=\alph*)]
\item Determine o domínio e a imagem de $f(x,y)=\sqrt{9-x^2-y^2}$.
\item Determine o domínio e a imagem do campo $\vet{F}(x,y)=(x+y,\,x-y)$.
\end{enumerate}
\end{exercicio}

\begin{exercicio}[{[Básico --- curvas de nível]}]
Descreva a curva de nível $c$ de $f(x,y)=x^2+y^2$ e a de $g(x,y)=x^2-y^2$, nos casos $c>0$, $c=0$ e $c<0$. Diga quais são vazias.
\end{exercicio}

\begin{exercicio}[{[Básico --- aberto, fechado, compacto]}]
Para cada conjunto, decida se é aberto, se é fechado, se é limitado e se é compacto. Justifique com a definição, não só com o nome.
\begin{enumerate}[label=\alph*)]
\item $A=\{(x,y):x^2+y^2<1\}$.
\item $B=\{(x,y):x^2+y^2\le 1\}$.
\item $C=\{(x,y):xy\ge 1\}$.
\item $D=\{(x,y):0<x\le 1\}$.
\end{enumerate}
\end{exercicio}

\begin{exercicio}[{[Básico --- $\eps$ e $\delta$]}]
Prove, pela definição, que
\[
\lim_{(x,y)\to(1,2)}(3x-y)=1.
\]
Exiba um $\delta$ em função de $\eps$.
\end{exercicio}

\begin{exercicio}[{[Intermediário --- cota pela distância]}]
Prove que
\[
\lim_{(x,y)\to(0,0)}\frac{x^3+y^3}{x^2+y^2}=0.
\]
A cota tem de valer no disco furado inteiro. Um caminho só não serve.
\end{exercicio}

\begin{exercicio}[{[Intermediário --- dois caminhos]}]
Mostre que
\[
\lim_{(x,y)\to(0,0)}\frac{x^2-y^2}{x^2+y^2}
\]
não existe, calculando o limite ao longo dos dois eixos.
\end{exercicio}

\begin{exercicio}[{[Intermediário --- polares]}]
Use $x=r\cos\theta$, $y=r\sin\theta$ para provar que
\[
\lim_{(x,y)\to(0,0)}\frac{\sin(x^2+y^2)}{x^2+y^2}=1.
\]
O $\delta$ tem de ser independente de $\theta$.
\end{exercicio}

\begin{exercicio}[{[Avançado --- retas não bastam]}]
Seja
\[
f(x,y)=\frac{x^2 y}{x^4+y^2},\qquad (x,y)\neq(0,0).
\]
\begin{enumerate}[label=\alph*)]
\item Calcule o limite de $f$ ao longo de cada reta $y=mx$ e ao longo de $x=0$.
\item Calcule o limite ao longo da parábola $y=x^2$.
\item Conclua, citando o teorema que você usa, que $\lim_{(x,y)\to(0,0)} f(x,y)$ não existe.
\end{enumerate}
\end{exercicio}

\begin{exercicio}[{[Avançado --- produto de limites]}]
Prove que, se $\lim f=L$ e $\lim g=M$ no mesmo ponto de $\R^2$, então $\lim(fg)=LM$. Escreva os quantificadores. O passo em que $f$ fica limitada perto do ponto precisa aparecer na conta, com a constante explícita.
\end{exercicio}

\end{document}
